\subsection{Модуль \texttt{families}}
\label{subsec:families-module}

Компонент \texttt{families} обеспечивает описание параметрических семейств распределений и построение конкретных экземпляров распределений из параметров. В отличие от \texttt{distributions}, который фиксирует «как вычислять», модуль \texttt{families} фиксирует «что такое распределение данного семейства»: набор параметров, допустимость значений, носитель и базовые (аналитические) характеристики.

\begin{figure}[H]
    \centering
    \includegraphics[width=.9\linewidth]{families.png}
    \caption{Ключевые классы модуля \texttt{families}: семейство, параметризации, распределение и реестр семейств}
    \label{fig:uml-families}
\end{figure}

\subsubsection{Семейство распределений: \texttt{ParametricFamily}}

\texttt{ParametricFamily} представляет одно семейство распределений и хранит:
\begin{itemize}[compact]
    \item набор доступных параметризаций и имя базовой параметризации;
    \item правила приведения параметров к базовой форме;
    \item аналитические реализации характеристик, привязанные к параметризациям;
    \item \texttt{support\_resolver} --- функцию/объект, выводящий носитель распределения по параметрам.
\end{itemize}

Семейство предоставляет фабричный метод \texttt{distribution(...)}. Он принимает значения параметров (и, опционально, имя параметризации), валидирует их и возвращает объект \texttt{ParametricFamilyDistribution}, реализующий протокол \texttt{Distribution}. Таким образом, дальнейшие вычисления выполняются уже в инфраструктуре \texttt{distributions}.

\subsubsection{Параметризации и ограничения}

Параметризация представлена классом \texttt{Parametrization}. Экземпляр параметризации содержит словарь значений параметров и предоставляет:
\begin{itemize}[compact]
    \item \texttt{validate()} --- проверку набора ограничений;
    \item \texttt{transform\_to\_base\_parametrization()} --- преобразование параметров к базовой параметризации семейства.
\end{itemize}

Ограничения оформляются объектами \texttt{ParametrizationConstraint} (описание + предикат \texttt{check}). Важное свойство такой схемы состоит в том, что ограничения задаются рядом с параметризацией и остаются локальными: добавление новой параметризации не требует модификации существующих.

\subsubsection{Аналитические характеристики и план вывода}

Семейство хранит набор аналитических реализаций характеристик (на уровне \texttt{GenericCharacteristicName}) и допускает, что разные характеристики удобнее задавать в разных параметризациях. Для этого используется \texttt{analytical\_plan}: он задаёт, к какой параметризации следует привести параметры, чтобы использовать доступную аналитическую формулу.

На уровне экземпляра распределения (\texttt{ParametricFamilyDistribution}) формируется отображение \texttt{analytical\_computations}. Эти вычисления становятся «якорями» для дальнейшего вывода характеристик по реестру (см.~\S\ref{subsec:characteristic-registry}).

\subsubsection{Экземпляр распределения: \texttt{ParametricFamilyDistribution}}

\texttt{ParametricFamilyDistribution} связывает воедино:
\begin{itemize}[compact]
    \item семейство и выбранную параметризацию;
    \item конкретные значения параметров (в том числе в базовой форме);
    \item \texttt{support} и \texttt{distribution\_type};
    \item стратегии вычислений и сэмплирования.
\end{itemize}

Такое распределение является полноценным объектом \texttt{Distribution}: оно может генерировать выборку через \texttt{sampling\_strategy} и вычислять характеристики через \texttt{computation\_strategy}.

\subsubsection{Глобальный реестр семейств}

Для поддержки готовых семейств и интеграции с другими пакетами экосистемы введён синглтон \texttt{ParametricFamilyRegister}. Он предоставляет операции регистрации семейства, получения по имени и перечисления зарегистрированных семейств. Наличие отдельного реестра позволяет подключать «поставки» семейств из разных модулей, не связывая их напрямую.
